\documentclass[11pt,a4paper]{article}
\usepackage[margin=1in]{geometry}
\usepackage{amsmath,amssymb}
\usepackage[numbers,sort&compress]{natbib}
\usepackage{hyperref}
\usepackage{graphicx}
\usepackage{bm}

\title{Truncated Marginal Neural Ratio Estimation (TMNRE)\\\large A Simulation-Based Inference Approach}
\author{Murali Krishnan}
\date{\today}

\begin{document}

\maketitle

\begin{abstract}
This report presents a structured overview of Truncated Marginal Neural Ratio Estimation (TMNRE), a neural simulation-based inference algorithm designed for efficient and credible Bayesian parameter estimation in settings with high-dimensional parameters and intractable likelihoods. The focus is on the core ideas, algorithmic details, and practical implementation aspects.
\end{abstract}

\section{Introduction}

Many physical models are most naturally specified as stochastic forward simulators: given parameters $\bm{\theta}$, one can generate synthetic data $\bm{x}$ by running a chain of physical processes, instrument responses, and noise realizations. Cosmological analyses of the cosmic microwave background, strong gravitational lensing images, type Ia supernova light curves, stochastic gravitational-wave background searches, and 21-cm cosmology all fall into this category~\citep{tmnreCosmology,tmnreStrongLensing,tmnreSupernovaeDust,tmnreGravitationalWaves,tmnre21cm}. In such problems the forward model is available, but the likelihood $p(\bm{x} \mid \bm{\theta})$---the probability of the observed data given parameters---is typically not, because writing it down would require integrating over an enormous number of latent variables introduced along the simulation chain~\citep{tmnreOriginal,tmnreGravitationalWaves}.

Simulation-based inference (SBI) is the family of methods that performs Bayesian inference in exactly this regime, using only samples $(\bm{\theta}, \bm{x}) \sim p(\bm{\theta})p(\bm{x} \mid \bm{\theta})$ drawn from the prior and the simulator~\citep{tmnreOriginal,autoregressiveNre}. The practical difficulty is cost. Classical approaches such as Approximate Bayesian Computation, or likelihood-based Markov Chain Monte Carlo (MCMC) with hand-crafted approximate likelihoods, scale poorly when the simulator is expensive and when the posterior occupies a tiny fraction of the prior volume~\citep{tmnreOriginal,tmnreCosmology}.

\subsection*{Two obstacles: dimensionality and trust}

The first obstacle is dimensionality, and it is worth being precise about why it bites. Suppose the prior is a product of independent distributions over $d$ parameters, and that for each parameter the posterior is narrower than the prior by a factor $f < 1$. The posterior then occupies a fraction $\sim f^{d}$ of the prior volume. For a modest $f = 0.1$ and $d = 20$, this is $10^{-20}$: a naive rejection or importance-sampling scheme that draws from the prior would need astronomically many simulator calls before a single sample landed in the region of interest. This exponential shrinkage of the useful fraction of parameter space is the sense in which SBI suffers from the curse of dimensionality~\citep{tmnreOriginal,tmnreCosmology}. The same exponential scaling afflicts attempts to model the full joint posterior density $p(\bm{\theta} \mid \bm{x}_0)$: the number of samples needed to constrain a $d$-dimensional density grows rapidly with $d$, even though in practice a physicist usually wants only a handful of one- and two-dimensional marginals for the parameters of interest, with the remaining nuisance parameters integrated out~\citep{mnreOverview,tmnrePerturberPopulation}.

The second obstacle is trust. Because the ground-truth posterior is inaccessible in a real analysis, one would like to test whether the reported credible intervals are actually calibrated---whether a nominal $68\%$ interval contains the true parameter $68\%$ of the time. Such tests require evaluating the inference procedure repeatedly on many mock observations, which is prohibitive for a sampling method that must be re-run from scratch for every dataset~\citep{tmnreOriginal,tmnreTalkCole}.

\subsection*{How TMNRE works}

Truncated Marginal Neural Ratio Estimation addresses both obstacles by combining three ingredients~\citep{tmnreOriginal,tmnreTalkWeniger}.

\paragraph{Ratio estimation instead of density estimation.} Rather than modelling a likelihood or posterior density directly, TMNRE trains a neural network as a binary classifier that distinguishes pairs drawn jointly, $(\bm{\theta}, \bm{x}) \sim p(\bm{\theta}, \bm{x})$, from pairs drawn independently, $(\bm{\theta}, \bm{x}) \sim p(\bm{\theta})p(\bm{x})$. At the optimum, the classifier's output recovers the likelihood-to-evidence ratio
\begin{equation}
  r(\bm{x} \mid \bm{\theta})
  = \frac{p(\bm{x} \mid \bm{\theta})}{p(\bm{x})}
  = \frac{p(\bm{\theta} \mid \bm{x})}{p(\bm{\theta})},
  \label{eq:ratio}
\end{equation}
so a trained network converts prior samples into weighted posterior samples without any likelihood evaluation~\citep{tmnreOriginal,tmnreExperimentsCode}.

\paragraph{Marginal rather than joint targets.} TMNRE trains a separate ratio estimator for each low-dimensional marginal of interest, typically the one- and two-dimensional slices of $\bm{\theta}$ that will appear in the final corner plot. Nuisance parameters are marginalized implicitly: they are varied in the simulations but simply not conditioned on in Eq.~\eqref{eq:ratio}. This sidesteps the high-dimensional density-estimation problem described above and, in the CMB application, made the number of simulations required effectively independent of the number of nuisance parameters~\citep{tmnreCosmology,mnreOverview,tmnreStrongLensing}.

\paragraph{Iterative truncation of the prior.} TMNRE proceeds in rounds. After each round of training, the learned marginals for the specific observation $\bm{x}_0$ are used to construct a constrained region $\Gamma \subset \Theta$ containing essentially all of the posterior mass, and the prior is multiplied by an indicator function on $\Gamma$~\citep{tmnreOriginal,swyftTruncationTutorial}. Concretely, for each parameter the algorithm discards the parameter values whose estimated marginal posterior density falls below a small threshold---the low-density regions far from the bulk of the posterior, which in a unimodal case are the outer wings of the marginal distribution. Inside $\Gamma$ the prior retains its original shape and normalization up to a constant; only the excluded outer regions change, and they are assigned zero mass. Simulations in the next round are drawn from this truncated prior, so simulator effort is concentrated where the posterior actually lives. The rounds terminate once successive truncated regions cease to shrink appreciably~\citep{tmnreOriginal,swyftTruncationTutorial}.

\subsection*{What this buys}

The gains from truncation are substantial and have been quantified in several analyses. In the CMB power-spectrum study of Cole et al., TMNRE recovered accurate one- and two-dimensional marginal posteriors for the six $\Lambda$CDM parameters from roughly $3 \times 10^{3}$ simulations, whereas the reference MCMC required $\mathcal{O}(5 \times 10^{6})$ simulator calls for converged $1$- and $2\sigma$ contours---about three orders of magnitude more~\citep{tmnreCosmology,coleSbiLectures}. For the stochastic gravitational-wave background, TMNRE needed $98\%$ fewer waveform evaluations than the nested-sampling approach in standard use~\citep{tmnreGravitationalWaves}. Reported improvements elsewhere are more modest but still significant, with factors of a few to a few tens in 21-cm and large-scale-structure applications, so the headline ``orders of magnitude'' claim should be read as problem-dependent rather than universal~\citep{tmnre21cm,mnreOverview}.

Separately from this efficiency gain, TMNRE is \emph{locally amortized}: the trained networks are valid for any data resembling $\bm{x}_0$, rather than only for $\bm{x}_0$ itself, though not for arbitrary datasets across the whole prior. Because evaluating a trained estimator is cheap, one can run it on many mock observations generated within the truncated region and thereby perform the coverage and consistency tests that motivated the second obstacle above---diagnostics that are impractical for a single MCMC chain~\citep{tmnreOriginal,tmnreTalkCole,autoregressiveNre}.

\subsection*{Structure of this report}

Section~\ref{sec:background} reviews the SBI setting more carefully. Section~\ref{sec:nre} develops neural and marginal ratio estimation. Section~\ref{sec:core} states the core TMNRE construction, and Section~\ref{sec:algorithm} gives the algorithm in full, including truncation criteria and stopping rules. Section~\ref{sec:implementation} covers implementation with \texttt{swyft}, the reference implementation of TMNRE~\citep{swyftSoftware,tmnreExperimentsCode}. Section~\ref{sec:examples} works through illustrative examples, Section~\ref{sec:diagnostics} discusses diagnostics and calibration, Section~\ref{sec:comparison} compares TMNRE with alternative SBI algorithms, and Section~\ref{sec:limitations} collects limitations and open directions.

\section{Background: Simulation-Based Inference}
\label{sec:background}

\section{Neural and Marginal Ratio Estimation}
\label{sec:nre}

\section{Core Ideas of TMNRE}
\label{sec:core}

\section{TMNRE Algorithm}
\label{sec:algorithm}

\section{Implementation with \texttt{swyft}}
\label{sec:implementation}

\section{Illustrative Examples}
\label{sec:examples}

\section{Diagnostics and Reliability}
\label{sec:diagnostics}

\section{Relations to Other SBI Methods}
\label{sec:comparison}

\section{Limitations and Future Work}
\label{sec:limitations}

\section{Conclusion}
\label{sec:conclusion}

\bibliographystyle{unsrtnat}
\bibliography{references}

\end{document}
